% 第二章：相关工作

\chapter{相关工作}

\section{大语言模型}

大语言模型（Large Language Models，LLMs）是基于Transformer架构的深度学习模型，利用自注意力机制捕获序列内的复杂依赖关系，能够处理和生成类人文本。以GPT-3\cite{brown2020gpt3}、GPT-4\cite{achiam2023gpt4}、Llama 3\cite{dubey2024llama3}和Mistral\cite{jiang2023mistral}为代表的大语言模型在文本生成、机器翻译、对话系统等任务上展现出卓越性能，已在金融、医疗、教育等多个行业实现大规模应用。

大语言模型的训练通常包含预训练和微调两个阶段。

\subsection{预训练}

在预训练阶段，模型通过因果语言建模目标（Causal Language Modeling）\cite{vaswani2017attention}在大规模语料库上学习通用的语言表示，即最大化序列中下一个词元的条件概率：
\begin{equation}
L_{\text{causal}} = - \sum_{i=1}^{N} \log P(x_i \mid x_{<i}; \theta)
\end{equation}
其中$x_i$表示第$i$个词元，$x_{<i}$表示其前驱词元序列，$\theta$为模型参数。预训练过程的核心是自注意力机制，该机制通过计算查询、键和值矩阵之间的相似度来捕获词元间的关联：
\begin{equation}
\text{Attention}(Q, K, V) = \text{softmax}\!\left(\frac{QK^\top}{\sqrt{d_k}}\right)V
\end{equation}
其中$Q$、$K$、$V$分别为查询、键和值矩阵，$d_k$为键向量的维度。

\subsection{微调}

预训练完成后，模型通过在特定任务数据集上优化任务相关的损失函数（如交叉熵损失）进行微调\cite{hu2022lora}，以适应下游应用需求。此外，基于人类反馈的强化学习（RLHF）等对齐方法\cite{ouyang2022rlhf}通过引入源自人类评估的奖励信号，进一步使模型输出与人类偏好保持一致，广泛应用于对话系统等场景。

尽管大语言模型取得了显著进展，其在安全性方面仍存在严重隐患。本文重点关注后门攻击这一训练期威胁，旨在通过开发稳健的防御机制为大语言模型的安全部署提供保障。

\section{后门攻击}

后门攻击（Backdoor Attack）是一种隐蔽的训练期安全威胁，攻击者通过在训练数据中植入中毒样本或直接修改模型参数，使模型在接收含有特定触发器的输入时产生恶意输出，而在正常输入上保持与未受攻击模型相同的行为。随着大语言模型的广泛部署和对海量数据集的依赖，后门攻击已成为一个日益严峻的安全问题。根据触发器设计方式的不同，现有攻击方法分为以下三类。

\subsection{单触发器后门攻击}

单触发器后门攻击（Single-trigger Backdoor Attack）使用单一固定触发器（如特定词语序列或文本模式）操控模型输出。Kurita等人\cite{kurita2020weight}针对BERT模型提出了干净标签后门攻击，通过使用稀有触发词进行微调诱导攻击者选定的词元预测。BadEdit\cite{li2024badedit}则证明，单触发器后门同样可以通过模型编辑的方式植入——直接修改模型权重以建立触发器与目标行为之间的关联，无需污染训练数据。Qi等人\cite{qi2021hidden}使用预先指定的句法模板作为触发器，生成流畅的中毒输入并抵抗简单的词过滤防御。此类设计依赖显式词汇线索，在一定程度上可被简单的异常词检测防御识别。

\subsection{多触发器后门攻击}

多触发器后门攻击（Multi-trigger Backdoor Attack）将触发器分散至多个输入组件，仅当所有触发器同时出现时才激活恶意有效载荷，从而在提升隐蔽性的同时增强触发控制。Huang等人\cite{huang2023cba}提出的复合后门攻击（CBA）将触发器分散至不同的提示片段（如用户查询与系统指令），大幅降低了良性输入意外激活后门的概率。Hou等人\cite{hou2024dual}提出的双触发器文本后门攻击仅当特定句法模板和虚拟语气模式同时出现时才被激活，相比基于词元的触发器具有更强的隐蔽性。

\subsection{无触发器后门攻击}

无触发器后门攻击（Triggerless Backdoor Attack）不依赖显式触发器字符串，而是在看似正常的输入中嵌入隐蔽的激活条件。后门可以在语义层面被触发——特定的含义、意图或话题隐式充当触发器；也可以在句法层面被触发——特定的结构模式充当隐藏触发器。Yan等人\cite{yan2024vpi}针对指令微调的大语言模型提出了虚拟提示注入（VPI），通过语义触发器使模型表现得如同攻击者选定的虚拟提示被隐式追加，无需任何显式触发器字符串即可操控输出。Hao等人\cite{hao2024dtba}提出了分布式触发器后门攻击（DTBA），将场景级触发器分散于多轮对话之中，检测难度极高。此类攻击对传统触发器过滤防御方法构成了更严峻的挑战。

\section{后门防御}

后门防御方法主要分为检测（Detection）和净化（Purification）两大方向。检测方法旨在识别后门模型或含触发器的输入，但无法将受损模型恢复至可安全部署的状态。净化方法直接消除后门行为同时保留良性功能，更具实用价值。本文聚焦于后门净化，现有方法可分为以下四类。

\subsection{基于微调的防御}

基于微调的防御（Fine-tuning-based Defense）通过在良性数据上训练来抑制后门模型的恶意行为，同时保留干净性能。NAD（Neural Attention Distillation）\cite{li2021nad}采用基于蒸馏的微调策略，由干净教师模型引导后门学生模型在干净数据上匹配中间层的注意力表示。W2SDefense\cite{yang2022w2sdefense}通过训练小型干净教师模型并借助参数高效微调（如LoRA\cite{hu2022lora}）蒸馏至后门学生模型，在无需完整重训练的条件下实现高效的后门遗忘。然而，此类方法计算代价较高，对高级后门攻击效果有限，且依赖目标任务的干净数据集。

\subsection{基于剪枝的防御}

基于剪枝的防御（Pruning-based Defense）通过剪除或屏蔽编码后门行为的参数子集实现净化，代表方法包括Fine-Pruning\cite{liu2018finepruning}、PURE\cite{zhu2024pure}和Obliviate\cite{li2024obliviate}等。Fine-Pruning移除干净输入激活值较低的神经元（随后进行轻量微调）；PURE通过剪枝后门相关注意力头并归一化激活来抑制恶意行为；Obliviate针对PEFT场景中嵌入在LoRA适配器中的后门，通过干净探测统计检测可疑维度并选择性清零实现净化；Wanda\cite{sun2023wanda}作为轻量级净化基线，通过剪除激活感知重要性分数$|w| \cdot |a|$较低的权重实现净化。

\subsection{基于参数编辑的防御}

基于参数编辑的防御（Parameter-editing-based Defense）将后门视为嵌入在模型权重中的局部关联，通过直接修改小部分参数消除后门。ROME\cite{meng2022rome}通过对MLP层进行低秩更新重写局部关联，在防御方能够提供显式反事实对的情况下消除后门；MEMIT\cite{meng2023memit}将该范式推广至大规模多关联的高效重写。此类方法不引入推理开销，但假设触发的不当行为是可识别的且后门是充分局部化的，对未知或语义触发器效果有限。

\subsection{基于嵌入的防御}

基于嵌入的防御（Embedding-based Defense）在表示空间而非直接在模型权重中运作，核心直觉在于后门行为在嵌入空间中表现为异常方向。BEEAR\cite{zeng2024beear}通过对抗性操控嵌入表示来识别和消除指令微调大语言模型中的安全后门；CROW\cite{min2024crow}利用后门触发器引发异常逐层表示不稳定性的观察，在小型干净提示集上进行带内部一致性正则化器的轻量微调，无需触发器知识或干净参考模型。

然而，上述净化方法要么依赖不切实际的先验信息，要么缺乏可复用的净化机制，限制了其在BDaaS场景中的实用性。

\section{权重空间表示}

近期研究揭示了模型行为可以直接在权重空间中被捕获、表达和操控这一重要规律，为本文提出的后门净化方法奠定了理论基础。

\subsection{任务向量}

Ilharco等人\cite{ilharco2023taskv}将模型在微调前后的参数差值形式化为\textbf{任务向量}（Task Vector）：令$\tau = \theta_{\text{fine-tuned}} - \theta_{\text{base}}$，这些向量展现出近似线性的代数性质，可以被加减以组合行为，或取反以抑制特定能力。任务向量还支持类比迁移——令$\Delta_{A \to B} = \theta_B - \theta_A$，将其应用于另一个模型$\theta_C$可得：
\begin{equation}
\theta_D \approx \theta_C + \Delta_{A \to B}
\end{equation}
即使未在目标任务上进行训练，也能实现一定程度的能力迁移。这一发现表明，有意义的模型行为变化可以通过权重空间中的向量运算来捕获和迁移。

\subsection{模型合并}

Wortsman等人\cite{wortsman2022soups}提出的模型合并（Model Soups）框架证明，多个从同一基础模型独立微调得到的模型通常处于共享的低损失区域，对其权重取平均可以在不增加推理成本的情况下超越最佳单一检查点的性能，并提升鲁棒性和泛化能力。这一经验性发现表明，有意义的权重更新在参数空间中具有可组合性，模型行为可以通过轻量的权重运算被组合、迁移乃至部分消除。

此外，模型引导（Model Steering）的相关工作表明，高层行为可以通过紧凑的线性方向来控制。通过对比不同属性下的模型状态（如有毒vs.无害），可以提取出可作为行为控制器的方向向量，在推理时可用于执行诱导模型行为转变的表示编辑。

\subsection{对后门净化的启示}

上述研究共同揭示了一个关键洞见：模型行为可以被表达为权重空间中的向量，并可通过向量运算加以操控。如果良性微调行为可以被编码为任务向量并进行迁移，那么恶意的后门行为也应当可以被抽象为参数空间中的向量，并通过向量操作加以识别和消除。本文提出的ProtoPurify方法正是基于这一洞见，通过提取代表性的后门原型向量，并利用其在权重空间中抑制与之对齐的后门相关参数分量，实现通用后门净化。
